% ABSTRACT
Most text corpora carry an implicit graph structure that standard language-model
pretraining discards: Wikipedia articles hyperlink to related articles, source files
import one another, and papers cite prior work. \method treats such a corpus as a
\emph{text-attributed graph}---documents are nodes, links are edges---and makes that
structure a first-class part of both pretraining and inference. Two ingredients: (i) a
graph traversal packs topologically close documents into each $32$k-token training
sequence, and (ii) a block-sparse attention mask keeps documents causally isolated
except along explicit edges, granting a linking document read-access, from the link
position onward, into the document it links to. The same link machinery runs at
inference: a generated link that resolves to a corpus document pulls that document
into the attention context under the training mask. On Wikipedia hyperlinks and code
imports across several languages, link-gated attention lowers the likelihood of
link-dependent completions relative to the same tokens scored without the target
($\Delta$nll $=\val{crossdoc.wiki_hotpotqa.delta_nll}$ on a HotpotQA cross-document
benchmark; \val{crossdoc.repobench_python.delta_nll}--\val{crossdoc.repobench_java.rw.delta_nll}
on RepoBench), leaves single-document benchmarks unchanged, and beats matched-compute
controls that grant more attention without the link gating. Holding the corpus fixed
and varying only the fraction of edges the mask may use, the benefit on code scales
near-linearly with realized link density, saturates early in training, and is
dominated by the density exposed at inference. A merged multi-source model
\fillin{strengthens/weakens?} the effect. We argue that the design's central promise
--- a model that fetches from its corpus by generating a reference, learned from
pretraining alone --- is a measurement these small models cannot yet support, and set
it up as the next step.
